\section{Глава 2. Многорукие бандиты}
\subsection{Упражнение 2.1(c. 54)}
Рассмотрим $\varepsilon$-жадный выбор действия для случая двух действий и $\varepsilon = 0.5$. Чему равна вероятность выбора жадного действия? 
\begin{solution}
    Для $\varepsilon$-жадной стратегии:
    \begin{itemize}
        \item с вероятностью $1-\varepsilon$ выбирается жадное действие.
        \item с вероятностью $\varepsilon$ выбирается случайное действие.
    \end{itemize}
    Поэтому $\mathbb{P}(\text{жадное}) = 1-\varepsilon+\varepsilon \cdot \frac{1}{2}=0.75$
\end{solution}

\subsection{Упражнение 2.2. Пример бандита(c. 57)}
Рассмотрим задачу о $k$-руком бандите с $k = 4$ действиями, обозначенными $1, 2, 3, 4$. Разберем применение к этой задаче алгоритма бандита с $\varepsilon$-жадным выбором действия, выборочным усреднением оценок ценности действий и начальными оценками $Q_1(a) = 0 \ \text{для} \ \forall a$. Предположим,
что начало последовательности действий и вознаграждений выглядит так: $A_1 = 1,
R_1 = 1, A_2 = 2, R_2 = -1, A_3 = 2, R_3 = -2, A_4 = 2, R_4 = 2, A_5 = 3, R_5 = 0$. На некоторых из этих временных шагов мог иметь место $\varepsilon$-случай, т.е. действие выбирается случайным образом. На каких шагах так оно заведомо и было? На каких шагах это могло случиться, а могло и не случиться?
\begin{solution} 
    Рассмотрим таблицу, где $Q_t(i), \ i=\overline{1,4}$ - значения оценок истинной ценности до выбора действия $A_t$:
    \renewcommand{\arraystretch}{1.3}
    \begin{table}[H]
    \centering
    \begin{tabular}{|c|c|c|c|c|c|c|c|}
    \hline
    $t$ 
    & $A_t$ 
    & $R_t$ 
    & $Q_t(1)$ 
    & $Q_t(2)$ 
    & $Q_t(3)$ 
    & $Q_t(4)$ 
    & $\varepsilon$-случай \\ \hline
    1 & 1 & 1 
    & 0 
    & 0 
    & 0 
    & 0 
    & мог быть / мог не быть \\
    \hline
    
    2 & 2 & $-1$ 
    & 1 
    & 0 
    & 0 
    & 0 
    & был заведомо \\
    \hline
    
    3 & 2 & $-2$ 
    & 1 
    & $-1$ 
    & 0 
    & 0 
    & был заведомо \\
    \hline
    
    4 & 2 & 2 
    & 1 
    & $-\frac{3}{2}$ 
    & 0 
    & 0 
    & был заведомо \\
    \hline
    
    5 & 3 & 0 
    & 1 
    & $-\tfrac{1}{3}$ 
    & 0 
    & 0 
    & был заведомо \\
    \hline
    
    6 & - & -
    & 1
    & $-\tfrac{1}{3}$ 
    & 0 
    & 0 
    & - \\
    \hline
    \end{tabular}
    \end{table}
\end{solution}

\subsection{Упражнение 2.3(с. 57)}
Глядя на рис. 2.2, можете ли вы сказать, какой метод будет работать лучше других в долгосрочной перспективе в терминах суммарного вознаграждения и вероятности выбора оптимального действия? Насколько лучше он будет? Выразите свой ответ количественно.
\begin{figure}
    \centering
    \includegraphics[width=1\linewidth]{fig_2_2.png}
    \caption{Рис. 2.2. Среднее качество $\varepsilon$-жадных методов оценки действий на 10-руком испытательном стенде. Данные получены усреднением по 2000 прогонов с разными
    задачами о бандите. Во всех методах использовалось выборочное усреднение оценок ценности действий}
\end{figure}
\begin{solution}
    \textbf{Вероятность.}
    Учтем, что при $\varepsilon > 0$ в долгосрочной перспективе оптимальное действие становится известно. Жадный алгоритм($\varepsilon = 0$) застревает в неверных шагах. Для вероятности выбора оптимального действия справедлива формула:
    \[\mathbb{P}(A_t=a^*) = (1 - \varepsilon) + \varepsilon \cdot \frac{1}{10}\]
    Тогда для
    \begin{itemize}
        \item $\varepsilon = 0.01: \quad \mathbb{P}(A_t=a^*) = 0.991$
        \item $\varepsilon = 0.1: \quad \mathbb{P}(A_t=a^*) = 0.91$
    \end{itemize}
    
    \textbf{Вознаграждение.} 
    Рассмотрим один временной шаг и запишем данные нам условия:
    \[q_*(1), \dots,q_*(10) \sim\mathcal N(0,1)\]
    \[R_t \mid (A_t = a) \sim\mathcal N(q_*(a),1)\]
    Используем закон полного математического ожидания:
    \[\mathbb{E}(R_t) = \mathbb{E}(\mathbb{E}(R_t \mid A_t)) = \mathbb{E}(q_*(A_t)),\]
    где ожидание берется и по $q_*$, и по случайности выбора действия $A_t$. \\
    Продолжая расписывать правую часть равенства, получаем
    \[\mathbb{E}(q_*(A_t)) = (1 - \varepsilon) \cdot \mathbb{E}(q_{max}) + 
    \varepsilon \cdot \mathbb{E}(\overline{q}),\]
    где $\mathbb{E}(\overline{q}) = \mathbb{E}_q(\frac{1}{10}\sum_{a=1}^{10}q_*(a)) = 0$ по условию задачи, а $q_{max} = \max_{a=1,\dots,10}q_*(a)$.\\
    В учебнике упоминается, что $\mathbb{E}(q_{max}) \approx 1.55$.\\
    Итого, мы получили:
    \[\mathbb{E}(R_t) \approx (1 - \varepsilon)\cdot1.55\]
    Тогда для
    \begin{itemize}
        \item $\varepsilon = 0.01: \quad \mathbb{E}(R_t) \approx 1.5345$
        \item $\varepsilon = 0.1: \quad \mathbb{E}(R_t) \approx 1.395$
    \end{itemize}
    Получается, что в среднем на одном временном шаге $\varepsilon = 0.01$ дает на $0.1395$ единиц вознаграждения больше, чем $\varepsilon = 0.1$.
\end{solution}

\subsection{Упражнение 2.4(с. 60)}
Если размеры шагов $\alpha_n$ непостоянны, то оценка $Q_n$ является взвешенным средним ранее полученных вознаграждений с весами, отличными от представленных в формуле (2.6). По аналогии с (2.6) выразите вес каждого предшествующего вознаграждения в терминах последовательности размеров шагов в общем случае. 
\begin{equation}
Q_{n+1} = (1 - \alpha)^n \cdot Q_1 + \sum_{i=1}^{n}\alpha \cdot (1-\alpha)^{n-i} \cdot R_i
\tag{2.6}
\end{equation}
\begin{solution}
    Начнём с общего правила обновления:
\[
Q_{n+1}
= Q_n + \alpha_n (R_n - Q_n)
= (1-\alpha_n) Q_n + \alpha_n R_n .
\]

Подставим выражение для $Q_n$:
\[
Q_n
= (1-\alpha_{n-1}) Q_{n-1} + \alpha_{n-1} R_{n-1}.
\]

Тогда
\[
Q_{n+1}
= (1-\alpha_n)\bigl[(1-\alpha_{n-1}) Q_{n-1}
+ \alpha_{n-1} R_{n-1}\bigr]
+ \alpha_n R_n.
\]

Раскрывая скобки, получаем:
\[
Q_{n+1}
= (1-\alpha_n)(1-\alpha_{n-1}) Q_{n-1}
+ (1-\alpha_n)\alpha_{n-1} R_{n-1}
+ \alpha_n R_n.
\]

Продолжая подстановку рекурсивно, имеем:
\begin{equation}
Q_{n+1}
= \Bigl(\prod_{i=1}^{n} (1-\alpha_i)\Bigr) Q_1
+ \sum_{k=1}^{n}
\Bigl(\alpha_k \prod_{i=k+1}^{n} (1-\alpha_i)\Bigr) R_k
\tag{$*$}
\end{equation}


Следовательно, вес вознаграждения $R_k$ равен:
\[
w_k^{(n)}
= \alpha_k \prod_{i=k+1}^{n} (1-\alpha_i),
\qquad k=1,\dots,n.
\]
\end{solution}

\subsection{Упражнение 2.5(с. 60)}
Спланируйте и поставьте эксперимент, демонстрирующий трудности, с которыми сталкиваются методы выборочного среднего в нестационарных задачах. Используйте модифицированную версию
$10$-рукого испытательного стенда, в которой все $q^*(a)$ вначале равны, а затем вычисляются путем независимых случайных блужданий (скажем, путем прибавления нормально распределенного приращения со средним $0$ и стандартным отклонением $0.01$ ко всем $q^*(a)$ на каждом шаге). Постройте графики, как на рис.2.2, для метода ценности действий, в  котором используются инкрементно вычисляемые выборочные средние, и еще одного метода, в котором размер шага постоянный, $\alpha = 0.1$. Положите $\varepsilon = 0.1$, а количество шагов возьмите большим, скажем $10 000$.
\begin{solution}
    \textbf{Код смотри в отдельном файле репозитория.}
    \begin{figure}[H]
        \centering
        \includegraphics[width=1\linewidth]{fig_2_5_sol_1.png}
    \end{figure}
    \begin{figure}[H]
        \centering
        \includegraphics[width=1\linewidth]{fig_2_5_sol_2.png}
    \end{figure}
\end{solution}
\notebook{02_05}{2.5}

\subsection{Упражнение 2.6. Таинственные пики(с. 62)}
Результаты, показанные на рис.2.3, должны быть вполне надежными, поскольку получены усреднением по 2000 независимых, случайно выбранных задач о 10-руком бандите. Но тогда откуда берутся осцилляции и  пики в  начальной части кривой для оптимистического метода? Иными словами, что могло бы заставить этот метод работать существенно хуже или лучше в среднем именно на первых шагах?
\begin{figure}
    \centering
    \includegraphics[width=1\linewidth]{fig_2_3.png}
    \caption{Рис. 2.3. Демонстрация оптимистических начальных оценок ценности действий на 10-руком испытательном стенде. В обоих методах размер шага постоянный, $\varepsilon$ = 0.1}
\end{figure}
\begin{solution}
    В первые шаги оптимистическая инициализация заставляет алгоритм почти детерминированно перебрать все действия, понижая их оценки. После первичного перебора оптимальное действие, имеющее более высокое истинное среднее, сохраняет более высокую оценку по сравнению с другими, что приводит к временному увеличению доли его выбора. Именно эта синхронная фаза перебора и последующее перераспределение оценок создают наблюдаемые пики и осцилляции.
\end{solution}

\subsection{Упражнение 2.7. Несмещенный постоянный размер шага(с. 62)}
В  большей части этой главы мы использовали выборочные средние для оценки ценности действий, поскольку при этом отсутствует начальное смещение, характерное для постоянного размера шага (см. анализ в формуле (2.6)). Однако выборочное среднее – не вполне удовлетворительное решение, поскольку оно может плохо работать в нестационарных задачах. Можно ли избежать смещения в  методах с  постоянным размером шага, не отказываясь от их преимуществ в нестационарных задачах?
Один из способов – использовать размер шага
\begin{equation}
\beta_n = \alpha / o_n
\tag{2.8}
\end{equation}
для обработки n-го вознаграждения за конкретное действие, где $\alpha > 0$ – обычный постоянный размер шага, а $o_n$
 – рекуррентная последовательность, начинающаяся в 0:
\begin{equation}
o_n = o_{n-1} + \alpha \cdot (1 - o_{n-1}) \quad \text{для} \quad n\ge0, \quad \text{где} \quad o_0=0
\tag{2.9}
\end{equation}
Проведите анализ, подобный формуле (2.6), и покажите, что $Q_n$
– экспоненциальное среднее, взвешенное по степени новизны, без начального смещения
\begin{solution}
    Заметим, что $o_n - \alpha = o_{n-1} \cdot (1 - \alpha)$.\\
    Согласно $(*)$ достаточно показать, что
    \[\prod_{i=1}^{n} (1-\beta_i) = 0\]
    Действительно, 
    \[\prod_{i=1}^{n} (1-\beta_i) = \prod_{i=1}^{n} \frac{o_i-\alpha}{o_i} = (1-\alpha) \cdot \prod_{i=1}^{n} \frac{o_{i-1}}{o_i} = 0,\]
    так как $\frac{o_{0}}{o_1} = 0$
\end{solution}

\subsection{Упражнение 2.8. Пики ВДГ(с. 63)}
На рис.2.4 график, соответствующий методу ВДГ, имеет отчетливый пик на 11-м шаге. Почему? Ответ будет считаться полным, если он объясняет, почему вознаграждение увеличивается на 11-м шаге и почему уменьшается на последующих шагах. Указание: если c = 1, то пик не так сильно выражен.
\begin{figure}[H]
    \centering
    \includegraphics[width=1\linewidth]{fig_2_4.png}
    \caption{Рис. 2.4. Среднее вознаграждение за выбор ВДГ-действия на 10-руком испытательном стенде. Как видно, ВДГ, вообще говоря, работает лучше, чем $\varepsilon$-жадный выбор действия, за исключением первых $k$ шагов, когда он случайно выбирает из еще не опробованных действий}
\end{figure}
\begin{solution}
    В течение первых десяти шагов описанный алгоритм выберет по одному разу каждое из действий в случайном порядке. 11-ый шаг - это первый шаг, когда алгоритм будет выбирать действие на основании $Q(a)$. Это означает выбор действия, которое принесло наибольшую награду за первые 10 шагов.Таким образом объясняется резкое возрастание графика. После этого слагаемое
    $c \cdot \sqrt{\frac{\ln t}{N_t(a)}}$ не позволяет долго оставаться на одном действии, и алгоритм продолжает исследовать уже с учетом "недостоверности". Падение после 11-го шага объясняется постепенным уточнением ценности и последующим переключением на менее исследованные действия.
\end{solution}

\subsection{Упражнение 2.9(с. 64)}
Покажите, что в случае двух действий распределение softmax совпадает с  определяемым логистической, или сигмоидной, функцией, часто используемой в статистике и в искусственных нейронных сетях.
\begin{solution}
    Пусть у нас действия $a \in \{1, 2\}$ с логитами $H_1 \ \text{и} \ H_2$ соответственно. Тогда softmax задаёт вероятности
    \[
    \pi(1)=\frac{e^{H_1}}{e^{H_1}+e^{H_2}},
    \qquad
    \pi(2)=\frac{e^{H_2}}{e^{H_1}+e^{H_2}}.
    \]
    Преобразуем $\pi(1)$:
    \[
    \pi(1)
    = \frac{e^{H_1}}{e^{H_1}+e^{H_2}}
    = \frac{1}{1+\frac{e^{H_2}}{e^{H_1}}}
    = \frac{1}{1+e^{H_2-H_1}}
    = \frac{1}{1+e^{-(H_1-H_2)}}.
    \]
    
    Обозначим $\Delta = H_1 - H_2$. Тогда
    \[
    \pi(1)=\frac{1}{1+e^{-\Delta}}
    = \sigma(\Delta),
    \]
    где сигмоидная функция определяется как
    \[
    \sigma(x)=\frac{1}{1+e^{-x}}.
    \]
    
    Аналогично,
    \[
    \pi(2)
    = \frac{e^{H_2}}{e^{H_1}+e^{H_2}}
    = \frac{1}{1+e^{H_1-H_2}}
    = \frac{1}{1+e^{\Delta}}
    = \sigma(-\Delta)
    = 1-\sigma(\Delta).
    \]
\end{solution}

\subsection{Упражнение 2.10(с. 68)}
Пусть имеется задача о 2-руком бандите, в которой истинные ценности действий случайным образом меняются от шага к шагу. Точнее, предположим, что на любом временном шаге истинные ценности действий 1 и 2 равны соответственно 0.1 и 0.2 с вероятностью 0.5 (случай A) и 0.9 и 0.8 с вероятностью 0.5 (случай B). Если мы не можем сказать, какой случай имеет место на данном шаге, то каково максимально возможное математическое ожидание успеха и как следует себя вести, чтобы достичь его? Теперь предположим, что на каждом шагенам сообщают, какой случай, A или B, имеет место (хотя истинной ценности действий мы по-прежнему не знаем). Это задача ассоциативного поиска. Какого максимального математического ожидания успеха можно достичь в этой задаче и как следует себя вести для его достижения?
\begin{solution}
    Если игрок не знает, какой случай имеет место на данном шаге, то он использует одну и ту же стратегию вне зависимости от случая. Посчитаем средние истинные ценности:
    \[\mathbb{E}(q_*(1)) = 0.5 \cdot 0.1 + 0.5 \cdot 0.9 = 0.5\]
    \[\mathbb{E}(q_*(2)) = 0.5 \cdot 0.2 + 0.5 \cdot 0.8 = 0.5\]
    Тогда 
    \[\mathbb{E}(R_t) = \pi(1) \cdot 0.5 + (1 - \pi(1)) \cdot 0.5 = 0.5\]
    Поэтому среднее вознаграждение будет равняться 0.5 вне зависимости от политики игрока. \\
    Если же игроку известен случай на текущем шаге, то в таком случае он властен использовать разные политики в зависимости от случая. Оптимальным будет следующий выбор:
    \[\pi(1 \mid A) = 0 \quad \text{и} \quad \pi(2 \mid A) = 1\]
    \[\pi(1 \mid B) = 1 \quad \text{и} \quad \pi(2 \mid B) = 0\]
    Тогда 
    \[\mathbb{E}(R_t) = 0.5 \cdot 0.2 + 0.5 \cdot 0.9 = 0.55\]
\end{solution}

\subsection{Упражнение 2.11(с. 71)}
Создайте рисунок, аналогичный рис.2.6, только для нестационарного случая, описанного в упражнении 2.5. Включите $\varepsilon$-жадный алгоритм с постоянным размером шага и $\varepsilon = 0.1$. Количество шагов пусть будет равным 200 000, а в качестве показателя качества каждого алгоритма и значения параметра возьмите среднее вознаграждение за последние 100 000 шагов. 
\begin{figure}[H]
    \centering
    \includegraphics[width=1\linewidth]{fig_2_6.png}
    \caption{Рис. 2.6. Параметрическое исследование различных алгоритмов решения задачи о бандите, описанных в этой главе. Каждая точка – это среднее вознаграждение, полученное алгоритмом за 1000 шагов при определенном значении параметра}
\end{figure}
\begin{solution}
    \textbf{Код смотри в отдельном файле репозитория.}
    \begin{figure}[H]
        \centering
        \includegraphics[width=1\linewidth]{fig_2_11_sol.png}
    \end{figure}
\end{solution}
\notebook{02_11}{2.11}
